% ============================================================================
% LATEX PROJECT REPORT TEMPLATE
% ============================================================================
% Based on academic standards (ENSI format)
% Author: Copilot
% Date: 2025-05-03
% ============================================================================

\documentclass[12pt,a4paper,onecolumn,oneside]{report}

% ============================================================================
% PACKAGES
% ============================================================================

% Language and encoding
\usepackage[utf-8]{inputenc}
\usepackage[french,english]{babel}
\selectlanguage{english}

% Margins and spacing
\usepackage[left=2.5cm,right=2.5cm,top=2.5cm,bottom=2.5cm]{geometry}
\usepackage{setspace}
\onehalfspacing % 1.5 line spacing

% Graphics and colors
\usepackage{graphicx}
\usepackage[table,xcdraw]{xcolor}
\usepackage{tikz}

% Tables
\usepackage{array}
\usepackage{multirow}
\usepackage{booktabs}
\usepackage{tabularx}

% Code listings
\usepackage{listings}
\usepackage{xcolor}

% Math
\usepackage{amsmath}
\usepackage{amssymb}
\usepackage{mathtools}

% References and bibliography
\usepackage[backend=biber,style=ieee,sorting=nty]{biblatex}
\addbibresource{references.bib}

% Links and URLs
\usepackage{hyperref}
\hypersetup{
    colorlinks=true,
    linkcolor=blue,
    citecolor=blue,
    urlcolor=blue,
    bookmarksnumbered=true
}

% Captions
\usepackage{caption}
\captionsetup[figure]{labelfont=bold, textfont=normalfont, skip=10pt}
\captionsetup[table]{labelfont=bold, textfont=normalfont, skip=10pt}

% TOC customization
\usepackage{tocloft}
\setcounter{tocdepth}{3}
\setcounter{secnumdepth}{3}

% Footnotes
\usepackage{footmisc}

% Page headers/footers
\usepackage{fancyhdr}
\pagestyle{fancy}
\fancyhf{}
\fancyhead[R]{\thepage}
\fancyhead[L]{\leftmark}
\renewcommand{\headrulewidth}{0.4pt}

% Acronyms
\usepackage{acro}

% ============================================================================
% CUSTOM COMMANDS & SETTINGS
% ============================================================================

% Code listing style
\lstset{
    language=Python,
    basicstyle=\ttfamily\small,
    keywordstyle=\color{blue},
    commentstyle=\color{gray},
    stringstyle=\color{red},
    breaklines=true,
    captionpos=b,
    numbers=left,
    numberstyle=\tiny\color{gray},
    showstringspaces=false,
    tabsize=2,
    backgroundcolor=\color{lightgray!10}
}

% Custom colors
\definecolor{lightgray}{rgb}{0.95,0.95,0.95}
\definecolor{darkblue}{rgb}{0.0,0.0,0.5}

% Section formatting
\titleformat{\chapter}[display]
  {\normalfont\Large\bfseries}
  {Chapter \thechapter}
  {20pt}
  {\Large}

% Figure reference command
\newcommand{\figref}[1]{Figure \ref{fig:#1}}
\newcommand{\tabref}[1]{Table \ref{tab:#1}}
\newcommand{\secref}[1]{Section \ref{sec:#1}}

% ============================================================================
% DOCUMENT METADATA
% ============================================================================

\title{%
    \textbf{PROJECT REPORT}\\
    \textbf{Design and Development Project}\\
    \vspace{0.5cm}
    \Large{Your Project Title Here}
}

\author{%
    \textbf{Author 1 Name}\\
    \textbf{Author 2 Name} (if applicable)\\
    \vspace{0.3cm}
    \normalsize{\textit{Supervisor: Dr./Prof. Name}}\\
    \normalsize{\textit{School: National School of Computer Science (ENSI)}}
}

\date{Academic Year: 2024/2025}

% ============================================================================
% BEGIN DOCUMENT
% ============================================================================

\begin{document}

% ============================================================================
% FRONT MATTER
% ============================================================================

% Title page
\maketitle

% Custom title page (alternative, more formal)
\thispagestyle{empty}
\begin{center}
    \vspace*{1cm}
    \includegraphics[width=0.3\textwidth]{../images/ensi_logo.png}\\
    \vspace{1cm}
    
    \large{\textbf{Université de la Manouba}}\\
    \large{\textbf{École Nationale des Sciences de l'Informatique}}\\
    \vspace{2cm}
    
    \textbf{\Large{RAPPORT DE PROJET}}\\
    \textbf{\Large{Projet de Conception et de Développement}}\\
    \vspace{1cm}
    
    {\LARGE \textbf{Your Project Title}}\\
    \vspace{0.5cm}
    
    \normalsize{Subtitle or description}\\
    \vspace{2cm}
    
    \textbf{Auteurs:}\\
    Author Name 1\\
    Author Name 2\\
    \vspace{1cm}
    
    \textbf{Encadrant:}\\
    Dr./Prof. Supervisor Name\\
    \vspace{1cm}
    
    \textbf{Année Universitaire:} 2024/2025\\
    \vspace{2cm}
    
    \textit{Appréciations et signature de l'encadrant}\\
    \vspace{3cm}
    
    \line(1,0){150}
\end{center}

\clearpage

% Appreciation page
\thispagestyle{empty}
\vspace*{5cm}
\begin{center}
    \textbf{\Large{APPRECIATION AND SUPERVISOR'S SIGNATURE}}
\end{center}
\vspace{3cm}

\noindent Supervisor's signature: \hrulefill

\noindent Grade: \hrulefill

\noindent Date: \hrulefill

\clearpage

% Acknowledgements
\chapter*{Acknowledgements}
\addcontentsline{toc}{chapter}{Acknowledgements}

We would like to express our deepest gratitude to Dr./Prof. \textbf{Supervisor Name} for his/her invaluable guidance, constant support, and insightful feedback throughout the entire duration of this project.

We also thank the faculty and staff of the National School of Computer Science (ENSI) for providing us with a solid academic foundation and the necessary resources to carry out this project.

Special thanks to [colleagues, collaborators, experts] for their assistance and valuable suggestions.

Finally, we would like to dedicate this work to our families for their unconditional support and encouragement throughout this academic journey.

\clearpage

% Abstract
\begin{abstract}
    \noindent This project focuses on [brief description of the problem]. The main objective is to [state your main goal]. We propose a solution based on [your approach/methodology]. The system was implemented using [technologies/tools used]. 
    
    Key results show [quantitative findings and achievements]. The system demonstrates [performance metrics]. This work contributes to [significance/impact], offering a [value proposition].
    
    \textbf{Keywords:} keyword1, keyword2, keyword3, keyword4, keyword5
\end{abstract}

\clearpage

% Table of Contents
\tableofcontents

\clearpage

% List of Figures
\listoffigures

\clearpage

% List of Tables
\listoftables

\clearpage

% ============================================================================
% MAIN CONTENT
% ============================================================================

\setcounter{page}{1}
\pagenumbering{arabic}

% ============================================================================
% CHAPTER 1: INTRODUCTION
% ============================================================================

\chapter{Introduction}
\label{chap:intro}

\section{Problem Statement}
\label{sec:intro:problem}

Provide a clear description of the problem your project addresses. Start broad and then narrow down to your specific focus.

\section{Motivation and Significance}
\label{sec:intro:motivation}

Explain why this problem matters and what impact solving it would have.

\section{Project Objectives}
\label{sec:intro:objectives}

List the main objectives and goals of your project:
\begin{enumerate}
    \item Objective 1: Description
    \item Objective 2: Description
    \item Objective 3: Description
\end{enumerate}

\section{Project Scope}
\label{sec:intro:scope}

Define what is included and excluded from this project.

\section{Report Organization}
\label{sec:intro:organization}

This report is organized as follows:
\begin{itemize}
    \item \textbf{Chapter 1:} Introduction and project overview
    \item \textbf{Chapter 2:} Literature review and theoretical background
    \item \textbf{Chapter 3:} Requirements analysis and specifications
    \item \textbf{Chapter 4:} Design and architecture
    \item \textbf{Chapter 5:} Implementation and results
    \item \textbf{Chapter 6:} Conclusion and future work
\end{itemize}

% ============================================================================
% CHAPTER 2: LITERATURE REVIEW & THEORETICAL BACKGROUND
% ============================================================================

\chapter{Literature Review and Theoretical Background}
\label{chap:literature}

\section{Theoretical Study}
\label{sec:theory}

\subsection{Fundamental Concepts}
\label{sec:theory:concepts}

Explain the foundational concepts necessary to understand your solution.

\subsection{Related Technologies}
\label{sec:theory:tech}

Discuss relevant technologies, frameworks, or algorithms.

\subsubsection{Algorithm 1: Name}
Describe the algorithm, its approach, and mathematical foundations.

\begin{equation}
    f(x) = \sum_{i=0}^{n} a_i x^i
    \label{eq:polynomial}
\end{equation}

See Equation \ref{eq:polynomial} for an example formula.

\subsubsection{Algorithm 2: Name}
Continue with additional algorithms as needed.

\section{Study of Existing Solutions}
\label{sec:existing}

\subsection{Solution 1: System A}
\label{sec:existing:sysA}

\textbf{Purpose:} What does this system solve?

\textbf{Approach:} How does it work?

\textbf{Strengths:}
\begin{itemize}
    \item Strength 1
    \item Strength 2
\end{itemize}

\textbf{Limitations:}
\begin{itemize}
    \item Limitation 1
    \item Limitation 2
\end{itemize}

\subsection{Solution 2: System B}
\label{sec:existing:sysB}

Continue with additional existing solutions.

\section{Critique of Existing Solutions}
\label{sec:critique}

Identify gaps and limitations in existing approaches.

\section{Proposed Solution}
\label{sec:proposed}

Briefly describe your proposed solution and how it addresses the identified gaps.

\section{Work Methodology}
\label{sec:methodology}

Describe your development approach (Agile, Waterfall, Iterative, etc.) and major phases:

\begin{table}[h]
    \centering
    \caption{Project Timeline}
    \label{tab:timeline}
    \begin{tabular}{|l|l|l|l|l|}
        \hline
        \textbf{Phase} & \textbf{Start} & \textbf{End} & \textbf{Duration} & \textbf{Status} \\
        \hline
        Requirements & Week 1 & Week 2 & 2 weeks & \checkmark Done \\
        \hline
        Design & Week 3 & Week 4 & 2 weeks & \checkmark Done \\
        \hline
        Implementation & Week 5 & Week 10 & 6 weeks & \checkmark Done \\
        \hline
        Testing & Week 11 & Week 12 & 2 weeks & \checkmark Done \\
        \hline
        Documentation & Week 13 & Week 14 & 2 weeks & \checkmark Done \\
        \hline
    \end{tabular}
\end{table}

\section{Conclusion}
\label{sec:lit:conclusion}

Summarize the key insights from the literature review and set up the transition to requirements analysis.

% ============================================================================
% CHAPTER 3: REQUIREMENTS & SPECIFICATIONS
% ============================================================================

\chapter{Requirements Analysis and Specifications}
\label{chap:requirements}

\section{Requirements Analysis}
\label{sec:req:analysis}

\subsection{Stakeholders}
\label{sec:req:stakeholders}

Identify all parties involved:
\begin{itemize}
    \item End users
    \item Project supervisor
    \item Development team
    \item System administrators
\end{itemize}

\subsection{Functional Requirements}
\label{sec:req:functional}

List what the system must do:

\begin{table}[h]
    \centering
    \caption{Functional Requirements}
    \label{tab:fr}
    \begin{tabular}{|l|p{8cm}|}
        \hline
        \textbf{ID} & \textbf{Requirement} \\
        \hline
        FR1 & The system shall allow users to upload a file in LAZ format \\
        \hline
        FR2 & The system shall automatically build an octree representation \\
        \hline
        FR3 & The system shall store metadata in MongoDB \\
        \hline
        FR4 & The system shall provide a REST API for querying nodes \\
        \hline
    \end{tabular}
\end{table}

\subsection{Non-Functional Requirements}
\label{sec:req:nonfunctional}

List system characteristics:

\begin{table}[h]
    \centering
    \caption{Non-Functional Requirements}
    \label{tab:nfr}
    \begin{tabular}{|l|l|p{5cm}|}
        \hline
        \textbf{Category} & \textbf{ID} & \textbf{Requirement} \\
        \hline
        Performance & NFR1 & Response time < 100ms for API queries \\
        \hline
        Scalability & NFR2 & Support files up to 1 GB without overflow \\
        \hline
        Security & NFR3 & All data encrypted in transit and at rest \\
        \hline
        Reliability & NFR4 & System uptime > 99.5\% \\
        \hline
    \end{tabular}
\end{table}

\section{Specifications}
\label{sec:specifications}

Convert requirements into detailed specifications.

\section{Conclusion}
\label{sec:req:conclusion}

Summary of requirements and transition to design phase.

% ============================================================================
% CHAPTER 4: DESIGN & ARCHITECTURE
% ============================================================================

\chapter{Design and Architecture}
\label{chap:design}

\section{System Architecture}
\label{sec:design:architecture}

\subsection{Overall Architecture}
\label{sec:design:overall}

Describe the high-level system architecture.

\begin{figure}[h]
    \centering
    \begin{tikzpicture}[node distance=2cm]
        \node (client) [draw, rectangle, minimum width=2cm] {Client};
        \node (api) [draw, rectangle, right of=client, minimum width=2cm] {FastAPI};
        \node (minio) [draw, rectangle, right of=api, minimum width=2cm] {MinIO};
        \node (mongodb) [draw, rectangle, below of=api, minimum width=2cm] {MongoDB};
        \node (pdal) [draw, rectangle, below of=minio, minimum width=2cm] {PDAL};
        
        \draw [->] (client) -- (api);
        \draw [->] (api) -- (minio);
        \draw [->] (api) -- (mongodb);
        \draw [->] (api) -- (pdal);
    \end{tikzpicture}
    \caption{System Architecture Diagram}
    \label{fig:architecture}
\end{figure}

\subsection{Frontend Design}
\label{sec:design:frontend}

Describe the user interface and frontend components.

\subsection{Backend Design}
\label{sec:design:backend}

Describe the backend services, API endpoints, and data models.

\subsubsection{Database Schema}
\label{sec:design:schema}

\begin{lstlisting}[language=json, caption=Sample Database Schema, label=lst:schema]
{
  "_id": ObjectId,
  "dataset_id": "string",
  "node_id": "string",
  "depth": 0,
  "bbox": {
    "min_x": 0.0,
    "max_x": 100.0,
    "min_y": 0.0,
    "max_y": 100.0,
    "min_z": 0.0,
    "max_z": 50.0
  },
  "point_count": 10000,
  "is_leaf": true,
  "children": ["0", "1"],
  "created_at": "2025-05-03T11:00:00Z"
}
\end{lstlisting}

\section{Detailed Design}
\label{sec:design:detailed}

\subsection{Component Interactions}
\label{sec:design:interactions}

Describe how components interact.

\subsection{Algorithm Design}
\label{sec:design:algorithm}

\begin{lstlisting}[language=Python, caption=Octree Building Algorithm, label=lst:algorithm]
def build_octree(input_file, max_depth, point_threshold):
    """Build octree from point cloud file."""
    root_bbox = get_bounding_box(input_file)
    
    def process_node(file, bbox, depth):
        point_count = get_point_count(file)
        
        if should_split(depth, point_count):
            # Sample points
            sampled = sample_nth(file, step)
            remainder = get_remainder(file, step)
            
            # Create octants
            for octant in bbox.split_into_octants():
                octant_file = crop_to_bbox(remainder, octant)
                process_node(octant_file, octant, depth + 1)
        else:
            # Leaf node
            save_node(file, depth)
    
    process_node(input_file, root_bbox, 0)
\end{lstlisting}

\section{Conclusion}
\label{sec:design:conclusion}

Summary of design and transition to implementation.

% ============================================================================
% CHAPTER 5: IMPLEMENTATION & RESULTS
% ============================================================================

\chapter{Implementation and Results}
\label{chap:implementation}

\section{Development Environment}
\label{sec:impl:environment}

\subsection{Hardware Configuration}
\label{sec:impl:hardware}

\begin{table}[h]
    \centering
    \caption{Hardware Configuration}
    \label{tab:hardware}
    \begin{tabular}{|l|l|}
        \hline
        \textbf{Component} & \textbf{Specification} \\
        \hline
        CPU & Intel Core i7-10700K \\
        \hline
        RAM & 32 GB DDR4 \\
        \hline
        Storage & 1 TB SSD \\
        \hline
        GPU & NVIDIA RTX 3080 \\
        \hline
    \end{tabular}
\end{table}

\subsection{Software Environment}
\label{sec:impl:software}

\begin{table}[h]
    \centering
    \caption{Software Stack}
    \label{tab:software}
    \begin{tabular}{|l|l|}
        \hline
        \textbf{Component} & \textbf{Technology} \\
        \hline
        Language & Python 3.12 \\
        \hline
        Framework & FastAPI \\
        \hline
        Database & MongoDB 6.0 \\
        \hline
        Storage & MinIO (S3-compatible) \\
        \hline
        Processing & PDAL 2.6 \\
        \hline
        Containerization & Docker \& Docker Compose \\
        \hline
    \end{tabular}
\end{table}

\section{Data Preparation}
\label{sec:impl:data}

Describe data collection, preprocessing, and statistics.

\section{Core Functionality Implementation}
\label{sec:impl:functionality}

\subsection{Module 1: Octree Builder}
\label{sec:impl:octree}

Implementation of the main octree building module.

\subsection{Module 2: PDAL Processor}
\label{sec:impl:pdal}

Implementation of PDAL integration for point cloud processing.

\section{Results and Validation}
\label{sec:impl:results}

\subsection{Performance Metrics}
\label{sec:impl:metrics}

\begin{table}[h]
    \centering
    \caption{Performance Results}
    \label{tab:results}
    \begin{tabular}{|l|r|r|c|}
        \hline
        \textbf{Metric} & \textbf{Value} & \textbf{Target} & \textbf{Status} \\
        \hline
        File Processing Time & 0.5s & < 1.0s & \checkmark \\
        \hline
        Memory Usage Peak & 256 MB & < 500 MB & \checkmark \\
        \hline
        Node Creation Accuracy & 100\% & 99\% & \checkmark \\
        \hline
        Zero Duplication & Yes & Yes & \checkmark \\
        \hline
    \end{tabular}
\end{table}

\subsection{Comparative Analysis}
\label{sec:impl:comparison}

\begin{table}[h]
    \centering
    \caption{Comparison with Baselines}
    \label{tab:comparison}
    \begin{tabular}{|l|r|r|r|}
        \hline
        \textbf{Algorithm} & \textbf{Accuracy} & \textbf{Speed} & \textbf{Memory} \\
        \hline
        U-NET & 94\% & 2.1s & 512 MB \\
        \hline
        Fuzzy C-Means & 89\% & 1.5s & 256 MB \\
        \hline
        Our Approach & 96\% & 0.8s & 128 MB \\
        \hline
    \end{tabular}
\end{table}

\section{Testing and Validation}
\label{sec:impl:testing}

\subsection{Unit Testing}
\label{sec:impl:unit}

Test coverage: 87\%

\subsection{Integration Testing}
\label{sec:impl:integration}

End-to-end workflow validation results.

\section{User Interface}
\label{sec:impl:ui}

\begin{figure}[h]
    \centering
    \includegraphics[width=0.8\textwidth]{../images/dashboard.png}
    \caption{Main Dashboard Interface}
    \label{fig:dashboard}
\end{figure}

\section{Challenges and Solutions}
\label{sec:impl:challenges}

\subsection{Challenge 1: Memory Overflow}

\textbf{Problem:} Large files caused memory overflow.

\textbf{Solution:} Implemented streaming pipeline with PDAL.

\textbf{Result:} Reduced peak memory by 75\%.

\section{Conclusion}
\label{sec:impl:conclusion}

Summary of implementation achievements and results.

% ============================================================================
% CHAPTER 6: CONCLUSION & FUTURE WORK
% ============================================================================

\chapter{Conclusion and Future Work}
\label{chap:conclusion}

\section{Main Achievements}
\label{sec:conclusion:achievements}

Summarize what was accomplished:
\begin{itemize}
    \item Achievement 1
    \item Achievement 2
    \item Achievement 3
\end{itemize}

\section{Challenges Overcome}
\label{sec:conclusion:challenges}

Describe obstacles faced and how they were resolved.

\section{Limitations}
\label{sec:conclusion:limitations}

Identify current system limitations and constraints.

\section{Future Work}
\label{sec:conclusion:future}

\subsection{Short-term (1-3 months)}

\begin{enumerate}
    \item Implement batch processing for multiple files
    \item Add distributed processing support
    \item Create comprehensive API documentation
\end{enumerate}

\subsection{Medium-term (3-6 months)}

\begin{enumerate}
    \item Multi-machine scaling via Apache Spark
    \item Real-time streaming point cloud processing
    \item Web-based visualization interface
\end{enumerate}

\subsection{Long-term (6+ months)}

\begin{enumerate}
    \item Support for additional point cloud formats
    \item GPU-accelerated octree construction
    \item ML-based automatic threshold optimization
\end{enumerate}

\section{Personal Reflection}
\label{sec:conclusion:reflection}

Reflect on what you learned from this project and how it has shaped your understanding or career interests.

\section{Final Remarks}
\label{sec:conclusion:remarks}

Closing thoughts on the project's significance and impact.

% ============================================================================
% BACK MATTER
% ============================================================================

\clearpage

% Bibliography
\printbibliography[title={Bibliography and References}]

% ============================================================================
% APPENDICES
% ============================================================================

\appendix

\chapter{API Documentation}
\label{app:api}

\section{Endpoints}
\label{sec:app:endpoints}

\subsection{Upload File}

\begin{lstlisting}[language=bash]
POST /lidar/upload
Content-Type: multipart/form-data

Parameters:
  - file: The LAZ/LAS file to upload
  - dataset_name: Name for the dataset

Response:
  {
    "dataset_id": "abc123",
    "filename": "terrain.laz",
    "dataset_name": "test_dataset"
  }
\end{lstlisting}

\subsection{Process Dataset}

\begin{lstlisting}[language=bash]
POST /lidar/process/{dataset_id}
Content-Type: application/json

Body:
  {
    "max_depth": 8,
    "point_threshold": 10000
  }
\end{lstlisting}

\chapter{Database Schema Details}
\label{app:schema}

Complete database schema documentation.

\chapter{Installation Guide}
\label{app:installation}

\section{Prerequisites}

\begin{enumerate}
    \item Docker and Docker Compose
    \item Python 3.12+
    \item Git
\end{enumerate}

\section{Setup Instructions}

\begin{lstlisting}[language=bash]
# Clone repository
git clone <repository-url>
cd project

# Start services
docker compose up --build

# Access services
# API: http://localhost:8000
# MinIO: http://localhost:9001
# MongoDB: http://localhost:27017
\end{lstlisting}

\chapter{Code Listings}
\label{app:code}

\section{Core Module}

\begin{lstlisting}[language=Python, caption=Main Octree Builder Module, label=lst:main]
from app.services.octree_builder import OctreeBuilder
from app.core.minio_client import get_minio_client

def build_octree(dataset_id, max_depth=8, point_threshold=1000000):
    """Main function to build octree."""
    minio_client = get_minio_client()
    
    builder = OctreeBuilder(
        minio_client=minio_client,
        dataset_id=dataset_id,
        max_depth=max_depth,
        point_threshold=point_threshold
    )
    
    nodes = builder.build_octree(
        input_file=f"datasets/{dataset_id}/input.laz",
        input_in_minio=True
    )
    
    return builder.get_stats()
\end{lstlisting}

\chapter{Test Results}
\label{app:tests}

\section{Unit Test Coverage}

\begin{table}[h]
    \centering
    \caption{Test Coverage Report}
    \label{tab:testcoverage}
    \begin{tabular}{|l|r|}
        \hline
        \textbf{Module} & \textbf{Coverage} \\
        \hline
        PDALProcessor & 92\% \\
        \hline
        OctreeBuilder & 85\% \\
        \hline
        MinIO Client & 90\% \\
        \hline
        Total & 87\% \\
        \hline
    \end{tabular}
\end{table}

% ============================================================================
% END OF DOCUMENT
% ============================================================================

\end{document}
